
This study measured effective rank of pre-trained projection weights in Mistral-7B, a 7B-parameter Transformer. Results show effective ranks of 1554-1648 at 99\% variance threshold, representing ~40\% of model dimension (4096). These values exceed commonly-assumed compression thresholds by 6-$7\times$ and typical LoRA adaptation ranks by $50\times$.

The findings establish that projection weights maintain near full-rank structure at this scale. Operator entropy does not decrease with layer depth ($\beta$ = +0.001453, p = 0.072). Both lines of evidence indicate that the low-rank assumption does not hold for pre-trained weights at 7B scale.

This measurement clarifies LoRA's mechanism: effectiveness stems from exploiting low-dimensional structure of task-specific adaptations, not from compressing low-rank weights. The $50\times$ rank gap demonstrates these are independent properties.

For compression research, the results establish boundary conditions: post-hoc techniques assuming bounded-state representations (r\_eff < 256) face empirical constraints at 7B scale. The measured ranks approach near full-rank, redirecting research toward methods that accommodate or do not assume low-rank weight structure.

Future work should extend measurements across model scales, architectures, and to runtime attention patterns. Cross-scale analysis will reveal whether rank properties scale with dimension or plateau. Runtime analysis will determine whether input-dependent patterns exhibit different rank characteristics than learned weights.

The path forward requires systematic measurement. Which structural properties generalize? Which are scale-dependent? Can we predict from pre-training dynamics whether models will exhibit particular rank properties? These questions demand empirical investigation grounded in direct measurement rather than inference from indirect evidence.
